% !PRINTED 123
\weberpage{123}
%<<CONT
ist. Wenn nun $c$ ein nicht zu $A$ gehöriges Element von $P$ ist, so wird $x^{-1}cx$ nicht für alle Elemente $x$ gleich $c$ sein. Es werden sich aber gewisse Elemente $b, b', \ldots$ unter den $x$ finden, darunter gewiss das Einheitselement, die der Bedingung
\[ x^{-1}cx = c \]
genügen, und diese Elemente bilden eine Gruppe, weil aus
\weq{3}{b^{-1}cb = c,\quad b'^{-1}cb' = c}
folgt, dass auch
\[ (b'b)^{-1}c(b'b) = b^{-1}b'^{-1}cb'b = c \]
sein muss. Diese Gruppe bezeichnen wir mit $B$ und ihren Grad und Index mit $\mu, \varepsilon$, so dass
\weq{4}{n = \mu\varepsilon}
wird; $B$ ist ein echter Theiler von $P$, und daher ist $\varepsilon$ jedenfalls grösser als 1, weil sonst $c$ in $A$ enthalten wäre, gegen die Annahme.

Ist $g$ ein Element von $P$, das nicht in $B$ vorkommt, und $b$ ein beliebiges Element von $B$, so ist nach (3):
\weq{5}{g^{-1}b^{-1}cbg = g^{-1}cg,}
was von $c$ verschieden ist, weil $g$ nicht in $B$ vorkommt. Wenn man also $P$ in die Nebengruppen zerlegt:
\weq{6}{P = B + Bg_1 + Bg_2 + \cdots + Bg_{\varepsilon-1},}
so erhält man, wenn man in $x^{-1}cx$ für $x$ alle Elemente einer dieser Nebengruppen setzt, immer dasselbe Element $g^{-1}cg$, und man erhält also $\varepsilon$ und nicht mehr Elemente, und jedes gleich oft:
\weq{7}{c,\ g_1^{-1}cg_1,\ g_2^{-1}cg_2,\ \ldots,\ g_{\varepsilon-1}^{-1}cg_{\varepsilon-1};}
alle diese Elemente, deren Gesammtheit wir mit $C$ bezeichnen wollen, sind unter einander und von den Elementen $a$ verschieden. $C$ ist aber keine Gruppe, da das Einheitselement darunter nicht vorkommt.

Ist nun mit $A$ und $C$ die Gesammtheit der Elemente von $P$ noch nicht erschöpft, so nehme man ein in $A$ und $C$ nicht vorkommendes Element $c'$ und bilde nach derselben Vorschrift die Gruppe $B'$ vom Grade $\mu'$ aller Elemente, die der Bedingung
\[ b'^{-1}c'b' = c' \]
genügen, und das System $C'$:
\weq{8}{c',\ g_1'^{-1}c'g_1',\ g_2'^{-1}c'g_2',\ \ldots,\ g_{\varepsilon'-1}'^{-1}c'g_{\varepsilon'-1}',}
%CONT>>
